\documentclass[10pt,a4paper]{article}

\usepackage[T1]{fontenc}
\usepackage[utf8]{inputenc}
\usepackage[catalan]{babel}
\usepackage{lmodern}
\usepackage[a4paper,margin=2cm]{geometry}
\usepackage{xcolor}
\usepackage{booktabs}
\usepackage{tabularx}
\usepackage{longtable}
\usepackage{enumitem}
\usepackage{amsmath}
\usepackage{amssymb}
\usepackage{listings}
\usepackage{csquotes}
\DeclareQuoteStyle{catalan}
  {\guillemotleft}{\guillemotright}
  {\textquotedblleft}{\textquotedblright}
\usepackage{microtype}
\usepackage[unicode]{hyperref}

\hypersetup{
  colorlinks=true,
  linkcolor=black!70!blue,
  urlcolor=blue!60!black,
  pdftitle={El torneig dels mags: full de referencia},
  pdfauthor={Gerard Casacuberta, Diego Suarez, Mikolaj Jaworski}
}

\lstset{
  basicstyle=\ttfamily\small,
  breaklines=true,
  frame=single,
  columns=fullflexible,
  keepspaces=true,
  showstringspaces=false
}

\setlength{\parindent}{0pt}
\setlength{\parskip}{5pt}
\setlist{itemsep=2pt,topsep=4pt}
\renewcommand{\arraystretch}{1.15}

\title{\textbf{El torneig dels mags}\\[0.3em]
\large Full de referència: comunicació, missatges, ping i 2FA}
\author{Gerard Casacuberta Albert \and Juan Diego Suárez P. \and Mikolaj Jaworski}
\date{\today}

\begin{document}

\maketitle

Aquest full resumeix les dades que cal tenir a mà per explicar o depurar el projecte:
el ritme de comunicació, tots els missatges del servidor, el càlcul del ping i el
funcionament de la verificació en dos passos. Els noms de fitxer i de camp són els del
codi; les xifres són les constants reals.

\textbf{Les quatre respostes, en curt:}

\begin{itemize}[nosep]
  \item \textbf{Ritme}: sondeig cada \textbf{100~ms} (10~Hz); no hi ha tick de
        simulació al servidor. Cada casella de moviment dura 300~ms.
  \item \textbf{Missatges}: \textbf{3} esdeveniments de servidor a client
        (\texttt{move}, \texttt{duel}, \texttt{error}) i \textbf{5} accions HTTP de
        client a servidor.
  \item \textbf{Ping}: $\mathrm{RTT}=r-e$ i
        $\mathrm{offset}=s+\mathrm{RTT}/2-r$; 8 mostres, offset de la més ràpida i
        mitjana del RTT.
  \item \textbf{2FA}: TOTP de 6 dígits cada 30~s; QR i activació amb codi; el login
        amb 2FA intercanvia un repte de 5~min per un token de 2~h.
\end{itemize}

% ---------------------------------------------------------------------------
\section{Dues aplicacions, un sol servidor}

\begin{table}[ht]
\small
\begin{tabularx}{\textwidth}{@{}p{4.2cm}p{2.9cm}X@{}}
\toprule
\textbf{Peça} & \textbf{Carpeta} & \textbf{Tecnologia i responsabilitat} \\
\midrule
Mòdul de Gestió d'Usuaris & \texttt{server/} & Express 4, Eta 4, SQLite
(\texttt{better-sqlite3}), Argon2id, JWT, \texttt{multer}, \texttt{otplib} i Brevo.
Registre, login, 2FA, perfil, recuperació, baneig i trofeus. \\
Mòdul d'Interacció (client) & \texttt{client/} & Vite, PixiJS\ v8,
\texttt{pixi-tiledmap} i \texttt{axios}. Pantalles, mapa isomètric, animacions, HUD i
so. \\
Mòdul d'Interacció (servidor) & \texttt{server/GameRoom.js} i \texttt{server/game/} &
Sales en memòria, estat, BFS, moviment, disputes, millores, victòria i trofeus. \\
\bottomrule
\end{tabularx}
\end{table}

\begin{itemize}[nosep]
  \item \textbf{Un sol procés}: \texttt{server/index.js} munta les pàgines Eta, l'API
        d'usuaris i \texttt{/api/game}. No hi ha WebSockets ni Colyseus; el joc és una
        API REST amb sondeig.
  \item \textbf{Base de dades}: \texttt{server/db/data.db} per defecte;
        \texttt{DB\_PATH} la pot moure. L'esquema es crea sol a l'arrencada i el seed
        només s'importa si \texttt{users} és buida.
  \item \textbf{Client}: \texttt{window.\_\_SERVER\_URL\_\_} (generat a
        \texttt{/config.js}) o \texttt{http://<host>:2567} en desenvolupament;
        \texttt{API = BASE + /api}.
  \item \textbf{Ports}: servidor 2567; Vite 5173; nginx del client 8080 a Compose i
        10000 a Render.
  \item \textbf{Ordre d'arrencada}: primer el servidor, després el client; el login
        passa pel servidor.
\end{itemize}

% ---------------------------------------------------------------------------
\section{Ritme de comunicació (el \emph{tick})}

\textbf{Resposta curta: 100~ms, és a dir 10~Hz.} No és un \emph{tick} del servidor:
el servidor no executa un bucle de simulació a intervals fixos. El client fa sondeig
(\emph{polling}) i el servidor respon a cada petició i als seus temporitzadors.

\begin{table}[ht]
\small
\begin{tabularx}{\textwidth}{@{}p{5.3cm}p{2.4cm}X@{}}
\toprule
\textbf{Constant} & \textbf{Valor} & \textbf{On i per a què} \\
\midrule
\texttt{POLL\_MS} & 100~ms & \texttt{client/src/game/connection.js}. Objectiu entre
l'inici d'un sondeig i el següent; si el RTT supera 100~ms, el següent surt tot seguit. \\
\texttt{MS\_PER\_TILE} & 300~ms & Servidor i client. Durada de l'animació d'una
casella; el màxim és 3,33 caselles per segon. \\
\texttt{SAMPLES} & 8 & Client. Finestra de mostres de RTT i de desfasament de rellotge. \\
\texttt{SWEEP\_MS} & 2~s & \texttt{server/routes/game.js}. Recorregut de sales per
expulsar jugadors inactius. \\
\texttt{PLAYER\_TIMEOUT} & 15~s & Servidor. Sense cap sondeig durant aquest temps, el
jugador surt de la sala. \\
\texttt{CLICK\_ARBITRATION\_MS} & 100~ms & Servidor. Finestra base per resoldre una
casella disputada; s'amplia amb el RTT. \\
\texttt{MAX\_EVENTS} & 200 & Servidor. Màxim d'esdeveniments en cua per sessió. \\
\texttt{MAX\_PLAYERS} & 6 & Servidor. Places per sala (es pot canviar amb
\texttt{MAX\_PLAYERS}). \\
\texttt{MIN\_PLAYERS} & 1 o 2 & 1 en desenvolupament; 2 si
\texttt{COOKIE\_SECURE=production}. \\
\texttt{MAX\_MOVE\_RANGE} & 7 & Servidor. Abast màxim; s'inicia amb 3. \\
\bottomrule
\end{tabularx}
\end{table}

El bucle del client és seqüencial, no un \texttt{setInterval}:

\begin{lstlisting}
poll():
  sentAt = Date.now()
  data = GET /api/game/:roomId?since=cursor&rtt=rtt
  receivedAt = Date.now()
  process(data.events)
  cursor = data.cursor
  wait max(0, 100 - (Date.now() - sentAt))
\end{lstlisting}

Conseqüències pràctiques:

\begin{itemize}[nosep]
  \item Amb un RTT baix, unes 10 peticions per segon i un decalatge d'entre 0 i 100~ms
        (més el RTT) per veure un canvi d'estat.
  \item Amb un RTT alt, la freqüència baixa (sempre hi ha una petició en vol com a
        màxim); no se n'apilen diverses.
  \item El servidor no envia res pel seu compte: la resposta a la petició és l'únic
        canal de servidor a client.
\end{itemize}

% ---------------------------------------------------------------------------
\section{Missatges del servidor}

El servidor només genera \textbf{tres tipus d'esdeveniment} de joc. Tots viatgen dins
de \texttt{events} a la resposta de \texttt{GET /api/game/:roomId}; no hi ha
WebSockets ni notificacions push.

\textbf{Resum: 3 missatges de servidor a client (move, duel, error) i 5 accions HTTP de
client a servidor.}

\begin{table}[ht]
\small
\begin{tabularx}{\textwidth}{@{}p{1.5cm}X@{}}
\toprule
\textbf{Missatge} & \textbf{Contingut i efecte} \\
\midrule
\texttt{move} &
\texttt{\{id, path, startTime, msPerTile\}}. El servidor l'envia a tothom quan un
candidat comença a caminar (guanyador o perdedor d'una disputa), quan un perdedor
retrocedeix i quan es cancel·la una reclamació. \texttt{path} és una llista de
\texttt{\{x,y\}} que inclou la casella inicial; \texttt{msPerTile} és 300. El client
substitueix o inicia l'animació i, si és el seu propi eco, conserva la predicció local
quan el camí coincideix. \\
\texttt{duel} &
\texttt{\{winnerId, x, y, loserIds, startTime, duration\}}. S'envia quan una casella
disputada es resol amb almenys un perdedor. \texttt{x,y} és la casella del guanyador;
\texttt{startTime} és l'instant de l'impacte i \texttt{duration} és 900~ms. El client
mostra \enquote{¡Expelliarmus!} i reprodueix \texttt{spell.mp3} en aquest instant. \\
\texttt{error} &
\texttt{\{message\}}. Només per a errors de joc visibles: equip ja ocupat o clic que
arriba quan la disputa ja està resolta i la seva marca és anterior a la del guanyador.
El client l'escriu a l'element d'avís (\texttt{lobbyInfo}), visible al lobby. Els
moviments invàlids es descarten en silenci; no generen \texttt{error}. \\
\bottomrule
\end{tabularx}
\end{table}

\subsection*{Com es lliuren}

\begin{itemize}[nosep]
  \item Cada sessió té una cua \texttt{session.events} amb un comptador global
        \texttt{eventSequence}. Cada esdeveniment porta \texttt{seq}.
  \item El client envia \texttt{since=cursor}; el servidor retorna només els
        esdeveniments amb \texttt{seq > since}, el nou \texttt{cursor} i
        \texttt{serverTime}.
  \item La cua es limita a 200 elements. Si un client s'atura molt, perd els més
        antics; el snapshot complet de la resposta següent torna a l'estat correcte.
  \item El snapshot inclou \texttt{state}, \texttt{events}, \texttt{cursor} i
        \texttt{serverTime}; \texttt{blocked} (la graella de col·lisions) només arriba
        amb \texttt{since = 0} i, quan el cursor avança, deixa de rebre's.
\end{itemize}

\subsection*{Peticions del client}

\begin{table}[ht]
\small
\begin{tabularx}{\textwidth}{@{}p{6.2cm}p{1.1cm}X@{}}
\toprule
\textbf{Petició} & \textbf{Auth} & \textbf{Efecte} \\
\midrule
\texttt{POST /api/game/join} & Bearer & Obre o recupera una sala. Cos
\texttt{\{roomId?\}}; retorna \texttt{roomId}, \texttt{playerId}, estat, events,
cursor i \texttt{serverTime}. \\
\texttt{GET /api/game/:roomId} & Bearer & Paràmetres \texttt{since} i \texttt{rtt}.
Retorna el snapshot i actualitza \texttt{lastSeen} i \texttt{player.ping}. \\
\texttt{POST /api/game/:roomId/team} & Bearer & Cos \texttt{\{team\}} (1--6). Només al
lobby; un equip ocupat retorna \texttt{error}. \\
\texttt{POST /api/game/:roomId/move} & Bearer & Cos \texttt{\{x, y, clickedAt\}}.
Valida i inicia o descarta el moviment. \\
\texttt{POST /api/game/:roomId/leave} & Bearer o cos & Surt de la sala. El
\texttt{sendBeacon} de \texttt{pagehide} envia \texttt{token} al cos perquè no pot
posar capçaleres. \\
\bottomrule
\end{tabularx}
\end{table}

L'API de joc \textbf{mai no accepta la cookie}: només el Bearer o un token al cos
(pensat per al beacon de sortida). Un repte de 2FA (camp \texttt{scope}) tampoc no hi
entra.

% ---------------------------------------------------------------------------
\section{Ping i sincronització del rellotge}

Cada sondeig mesura el temps d'anada i tornada i el desfasament entre els dos
rellotges. Amb $\mathrm{RTT}_i$ la mostra $i$, $e_i$ l'instant d'enviament al client,
$r_i$ l'instant de recepció i $s_i$ el \texttt{serverTime} de la resposta:

\[
  \mathrm{RTT}_i = r_i - e_i,
  \qquad
  \mathrm{offset}_i = s_i + \frac{\mathrm{RTT}_i}{2} - r_i.
\]

\begin{enumerate}[nosep]
  \item El client guarda les \textbf{8 darreres} mostres.
  \item L'\texttt{offset} global és el de la mostra amb el \textbf{RTT mínim}, no la
        mitjana: la resposta més ràpida té menys error de cua i evita que l'animació
        tremoli.
  \item El RTT que s'envia al servidor és la \textbf{mitjana arrodonida} de les mostres
        disponibles (fins a vuit); viatja com a \texttt{rtt} al següent \texttt{GET}.
  \item El servidor valida el valor (0--5000~ms), l'arrodoneix i el desa a
        \texttt{player.ping}. El panell de ping mostra aquest valor, no una mesura
        local.
  \item El rellotge del joc és \texttt{serverNow() = Date.now() + offset}; amb ell
        s'inicien les animacions i es programa el so del duel.
\end{enumerate}

\textbf{Exemple}: si el client envia a $e=1000$~ms, rep a $r=1040$~ms i el servidor diu
$s=1020$~ms, aleshores $\mathrm{RTT}=40$~ms i
$\mathrm{offset}=1020+20-1040=0$~ms. Si el rellotge del client anés 15~ms avançat,
l'\texttt{offset} sortiria $-15$~ms i el client restaria aquest valor.

\subsection*{Finestra d'arbitratge}

Quan arriba el primer clic a una casella, el servidor obre una reclamació i espera:

\[
  W = \min\left(1500,\ \max\left(100,\ \left\lceil \frac{\mathrm{RTT}_{\max}}{2}
  \right\rceil + 50\right)\right)\ \text{ms},
\]

on $\mathrm{RTT}_{\max}$ és el ping més alt de la sala en aquell moment. Després ordena
els candidats per \texttt{clickedAt} i, si empaten, per ordre d'arribada. Un
\texttt{clickedAt} només s'accepta si és dins de $[\text{ara}-5000,\ \text{ara}+100]$
ms; si no, s'usa l'hora de recepció. El guanyador fa el camí sencer; els perdedors
avancen fins a la casella anterior i, quan el guanyador arriba, retrocedeixen un 30\%
dels passos (mínim un).

\subsection*{Colors del panell}

\begin{table}[ht]
\small
\begin{tabular}{@{}lll@{}}
\toprule
\textbf{Ping} & \textbf{Color} & \textbf{Significat} \\
\midrule
$< 0$ & gris clar & encara no hi ha cap mostra \\
$< 100$~ms & verd & bo \\
$< 250$~ms & groc & acceptable \\
$\geq 250$~ms & vermell & alt \\
\bottomrule
\end{tabular}
\end{table}

El panell s'obre amb el Tab o amb el botó \enquote{Ping}; només és visible durant
\texttt{playing}. El ping és informatiu: el servidor l'usa per ampliar la finestra
d'arbitratge, no per decidir el guanyador.

% ---------------------------------------------------------------------------
\section{Verificació en dos passos (2FA)}

És TOTP (RFC 6238): \textbf{6 dígits}, període de \textbf{30~s}, HMAC-SHA-1 i secret
Base32. L'emissor és \enquote{El torneo de los magos} i l'etiqueta és el nom d'usuari.
El rellotge del dispositiu ha d'estar sincronitzat; la verificació només accepta el
període actual.

\subsection*{Flux d'activació}

\begin{enumerate}[nosep]
  \item L'usuari, amb sessió, crida \texttt{POST /api/2fa/setup}. El servidor genera
        un secret, el desa com a pendent (\texttt{totp\_enabled = 0}) i retorna
        \texttt{\{secret, qr\}}, amb el QR com a imatge PNG incrustada.
  \item L'usuari escaneja el QR o introdueix la clau manual i envia un codi a
        \texttt{POST /api/2fa/enable}. Només si el codi és correcte s'activa
        \texttt{totp\_enabled = 1}.
  \item Per desactivar-lo, \texttt{POST /api/2fa/disable} exigeix un codi actual: una
        sessió robada no pot llevar el segon factor.
\end{enumerate}

\subsection*{Flux de login}

\begin{enumerate}[nosep]
  \item \texttt{POST /api/login} amb usuari i contrasenya. Si el compte té 2FA, el
        servidor \textbf{no} emet cap token de sessió: retorna
        \texttt{\{requires2fa: true, challenge\}}.
  \item El client mostra el camp de 6 dígits. \texttt{POST /api/login/2fa} amb
        \texttt{\{challenge, code\}} valida el repte i el codi i aleshores retorna el
        token de sessió de 2~h i posa la cookie.
  \item El repte és un JWT de \textbf{5 minuts} amb \texttt{scope} igual a
        \texttt{2fa}. Només prova que la contrasenya era correcta: no és una sessió.
        \texttt{attachUser} el rebutja i l'API de joc també.
\end{enumerate}

\begin{table}[ht]
\small
\begin{tabularx}{\textwidth}{@{}p{2.1cm}p{4.2cm}p{1.2cm}X@{}}
\toprule
\textbf{Mètode} & \textbf{Ruta} & \textbf{Auth} & \textbf{Què fa} \\
\midrule
POST & \texttt{/api/2fa/setup} & sessió & Genera secret pendent i QR; 409 si ja està
actiu. \\
POST & \texttt{/api/2fa/enable} & sessió & Activa amb un codi de 6 dígits; 400 si és
incorrecte; 429 si es fan massa intents. \\
POST & \texttt{/api/2fa/disable} & sessió & Desactiva amb un codi actual; esborra el
secret. \\
POST & \texttt{/api/login} & --- & Contrasenya; retorna el repte si cal 2FA o el token
si no. \\
POST & \texttt{/api/login/2fa} & repte & Repte + codi; retorna el token de 2~h. \\
\bottomrule
\end{tabularx}
\end{table}

\subsection*{Detalls que es pregunten sovint}

\begin{itemize}[nosep]
  \item \textbf{Límit}: 10 intents per usuari cada 15 minuts a \texttt{enable},
        \texttt{disable} i \texttt{login/2fa}. La clau és el nom d'usuari, no la IP.
  \item \textbf{Cache}: totes les rutes de 2FA responen
        \texttt{Cache-Control: no-store}; el QR i els codis no es poden cachejar pel
        navegador ni per un proxy.
  \item \textbf{Secret pendent}: \texttt{setup} el desa abans de confirmar-lo. Si
        l'usuari abandona, no protegeix el compte; un nou \texttt{setup} el
        substitueix.
  \item \textbf{Compte bloquejat}: el segon pas també comprova el baneig i retorna
        403 amb el motiu.
  \item \textbf{Base de dades}: \texttt{users.totp\_secret} (Base32) i
        \texttt{users.totp\_enabled} (0/1). No hi ha taula separada.
  \item \textbf{Codis de recuperació}: no n'hi ha. Si es perd l'aplicació, cal
        intervenció manual; és una limitació coneguda.
  \item \textbf{Sessions obertes}: activar o desactivar el 2FA no revoca els JWT que ja
        existeixen; duren 2~h.
\end{itemize}

% ---------------------------------------------------------------------------
\section{Sales, equips i moviment}

\begin{itemize}[nosep]
  \item \textbf{Codi de sala}: 6 caràcters de
        \texttt{ABCDEFGHJKLMNPQRSTUVWXYZ23456789} (sense I, O, 0 ni 1). Es genera amb
        \texttt{Math.random} i es comprova que no existeixi.
  \item \textbf{Entrada}: \texttt{join} sense codi reutilitza la sala on ja és el
        jugador o en crea una; amb codi, s'hi uneix. Màxim 6 jugadors; una sala
        \texttt{locked} o plena retorna 409.
  \item \textbf{Equips}: 1--6, únics. La partida comença quan tots els jugadors tenen
        equip i s'arriba a \texttt{MIN\_PLAYERS}. Els sis spawns surten de
        \texttt{spawnPositions}.
  \item \textbf{Moviment}: BFS de 4 direccions al servidor; comprova fase, jugador que
        no estigui caminant, límits, col·lisió, ocupació i
        \texttt{path.length - 1 <= min(moveRange, 7)}. Un moviment invàlid es descarta
        sense resposta d'error.
  \item \textbf{Abast}: 3 inicial; cada bola suma 1 fins a 7. Amb 7, passar per sobre
        d'una bola no la recull; queda per als altres.
  \item \textbf{Victòria}: arribar a la copa, també si es travessa en un recorregut.
        El servidor passa a \texttt{finished}, cancel·la reclamacions, desa el
        guanyador i incrementa \texttt{users.trophies} via \texttt{GameRoom.onWin}.
  \item \textbf{Sortida}: \texttt{leave} o beacon de \texttt{pagehide}. Si el navegador
        mor sense avisar, l'escombrador expulsa al cap de 15~s. En sortir l'últim, la
        sala s'esborra.
  \item \textbf{Revanxa}: \texttt{play-again} surt de la sala i fa un \texttt{join}
        nou; els altres poden entrar amb el mateix flux.
\end{itemize}

El snapshot de l'estat conté \texttt{phase}, \texttt{players}, \texttt{boosts},
\texttt{trophyX}, \texttt{trophyY}, \texttt{trophyAvailable}, \texttt{winnerId},
\texttt{winnerName} i \texttt{maxMoveRange}; \texttt{blocked} només arriba quan
\texttt{since=0}. \texttt{players} és un objecte indexat per
id amb \texttt{name}, \texttt{team}, \texttt{x}, \texttt{y}, \texttt{moving},
\texttt{moveRange} i \texttt{ping}.

% ---------------------------------------------------------------------------
\section{Autenticació, usuaris i recuperació}

\begin{table}[ht]
\small
\begin{tabularx}{\textwidth}{@{}p{5.2cm}X@{}}
\toprule
\textbf{Ruta} & \textbf{Resum} \\
\midrule
\texttt{POST /api/register} & Usuari (mín. 3), correu, naixement (edat mínima 16),
contrasenya (mín. 12) i confirmació opcional. Retorna 201 amb token i cookie; 409 si
l'usuari o el correu existeixen. \\
\texttt{POST /api/login} & Argon2; 403 si està banejat; \texttt{requires2fa} si cal;
si no, token de 2~h. \\
\texttt{POST /api/login/2fa} & Repte de 5~min + codi TOTP; retorna el token. \\
\texttt{POST /api/logout} & Esborra la cookie. \\
\texttt{GET /api/me} & Dades del compte, \texttt{totpEnabled} i
\texttt{trophies}. \\
\texttt{PUT /api/profile} & Sobrenom obligatori i avatar opcional (2~MB; JPEG, PNG o
WebP; nom UUID; esborra l'anterior). \\
\texttt{POST /api/password} & Contrasenya actual + nova + confirmació; política de 12
caràcters. \\
\texttt{POST /api/delete} & Contrasenya actual; esborra tokens i usuari en una
transacció; 204. \\
\texttt{POST /api/forgot-password} & Resposta genèrica; límits i enviament per Brevo. \\
\texttt{POST /api/reset-password} & Token + contrasenya nova; consum únic i
transaccional. \\
\bottomrule
\end{tabularx}
\end{table}

\begin{itemize}[nosep]
  \item \textbf{Sessió}: JWT de 2~h amb \texttt{id} i \texttt{username}; cookie
        \texttt{httpOnly}, \texttt{sameSite=lax}, 1~h de durada i \texttt{secure} amb
        \texttt{COOKIE\_SECURE=production}. El client Vite el guarda a
        \texttt{sessionStorage} i l'envia com a Bearer.
  \item \textbf{Contrasenya}: Argon2id; mínim 12 caràcters Unicode; sense exigir
        majúscules ni símbols; HIBP amb \emph{k-anonymity} (només els 5 primers
        caràcters del SHA-1). Si HIBP cau, es permet continuar i es registra
        \texttt{HIBP\_UNAVAILABLE}.
  \item \textbf{Recuperació}: 32 bytes aleatoris; només el SHA-256 a SQLite; 30~min;
        un sol ús; límits 5/IP/15~min per sol·licitar, 10/IP/15~min per restablir,
        1/correu/min i 3/correu/hora; enllaç
        \texttt{SERVER\_URL/reset\#reset-password=<token>}; Brevo REST.
  \item \textbf{Baneig}: el middleware global consulta la BD a cada petició
        autenticada. Un compte bloquejat només pot veure \texttt{/banned} i fer
        \texttt{/api/logout}; la resta de pàgines redirigeixen i l'API retorna 403 amb
        \texttt{reason} i \texttt{redirect}.
  \item \textbf{Pàgines Eta}: si \texttt{CLIENT\_URL} està definida, \texttt{/},
        \texttt{/login}, \texttt{/register}, \texttt{/profile}, \texttt{/forgot},
        \texttt{/reset} i \texttt{/delete} redirigeixen al client Vite; el fragment del
        reset es conserva.
\end{itemize}

% ---------------------------------------------------------------------------
\section{Mapa, render i client}

\begin{itemize}[nosep]
  \item \textbf{Mapa}: \texttt{client/public/assets/map.tmj}, $60 \times 60$ caselles de
        $232 \times 110$ píxels, orientació isomètrica. Capes \texttt{ground},
        \texttt{grass}, \texttt{sanctuary}, \texttt{walls}, \texttt{details} i
        \texttt{collision}.
  \item \textbf{Propietats}: \texttt{spawnPositions} (sis spawns),
        \texttt{trophyPosition} ($[55,55]$) i \texttt{boostPositions} (vuit boles).
  \item \textbf{Servidor}: llegeix el tmj una vegada en arrencar; en deriva
        \texttt{GRID} i la llista de caselles bloquejades. \textbf{Cal reiniciar el
        servidor després de canviar el mapa.} El client té \texttt{GRID = 60} fix a
        \texttt{utils/constants.js}.
  \item \textbf{Render}: \texttt{renderTerrainSurfaces} pinta \texttt{sanctuary},
        \texttt{grass} i \texttt{ground} en una sola passada, un sprite per casella, i
        amaga les capes originals perquè les cares laterals no tapin el terra veí.
  \item \textbf{Personatges}: sis fulls (una casa per equip); tres files (perfil,
        davant, darrere); fotogrames retallats pels píxels opacs i escalats a
        $88 \times 128$; direcció segons $dx$/$dy$; mirall per a l'esquerra.
  \item \textbf{Càmera}: escala base 0,6 multiplicada per
        $\max(\text{amplada}/1920, \text{alçada}/910)$ i seguiment del jugador local;
        la referència cobreix la finestra i no revela més mapa.
  \item \textbf{So}: \texttt{menu.mp3} i \texttt{game.mp3} amb crossfade;
        \texttt{spell.mp3} al duel, \texttt{win.mp3} a \texttt{finished} i
        \texttt{click.mp3} en canviar de pantalla; volums a \texttt{localStorage}.
  \item \textbf{Ajudes}: Tab obre el panell de ping; Esc obre o tanca els ajustos; el
        panell no és modal i no atura la partida.
\end{itemize}

% ---------------------------------------------------------------------------
\section{Desplegament i variables}

\begin{table}[ht]
\small
\begin{tabularx}{\textwidth}{@{}p{3.7cm}X@{}}
\toprule
\textbf{Variable} & \textbf{Efecte} \\
\midrule
\texttt{PORT} & Port del servidor (2567 per defecte). \\
\texttt{JWT\_SECRET} & Signatura dels JWT\@. Obligatòria a producció; en local hi ha un
valor de desenvolupament. \\
\texttt{COOKIE\_SECURE} & \texttt{production} fa la cookie \texttt{secure} i activa el
mínim de 2 jugadors. \\
\texttt{CLIENT\_ORIGIN} & Orígens permesos per CORS, separats per comes; amb
credencials. \\
\texttt{CLIENT\_URL} & URL del client Vite; redirigeix les pàgines Eta antigues. \\
\texttt{SERVER\_URL} & URL pública del servidor; s'usa als enllaços de recuperació.
A producció cal definir-la (o \texttt{RENDER\_EXTERNAL\_URL} a Render). \\
\texttt{DB\_PATH} & Fitxer SQLite; per defecte \texttt{server/db/data.db}; a producció
\texttt{/data/data.db}. \\
\texttt{UPLOAD\_DIR} & Carpeta dels avatars; per defecte
\texttt{server/public/uploads}; a producció \texttt{/data/uploads}. \\
\texttt{MAP\_PATH} & Ruta del tmj que llegeix el servidor. \\
\texttt{MAX\_PLAYERS} & Places per sala (6 per defecte). \\
\texttt{TRUST\_PROXY\_HOPS} & Nombre de proxys de confiança (0--5); 1 a Render. \\
\texttt{BREVO\_API\_KEY}, \texttt{BREVO\_SENDER\_EMAIL}, \texttt{BREVO\_SENDER\_NAME} &
Credencials i remitent del correu de recuperació. Només al servidor. \\
\bottomrule
\end{tabularx}
\end{table}

\begin{itemize}[nosep]
  \item \textbf{Local}: \texttt{npm start} o \texttt{npm run dev} creen \texttt{.env}
        a partir de \texttt{.env.example} si no existeix (mai a producció) i preparen
        l'esquema SQLite.
  \item \textbf{Docker}: \texttt{docker compose up --build}; servidor 2567 i client
        nginx 8080; volum \texttt{iso-data} per a \texttt{/data}. El context del
        servidor és l'arrel perquè copia \texttt{client/public/assets/map.tmj}.
  \item \textbf{Render}: dos serveis Docker; el servidor necessita el pla
        \texttt{starter} i un disc d'1~GB; el client genera \texttt{/config.js} amb
        \texttt{SERVER\_URL}.
  \item \textbf{Comprovació}: \texttt{GET /healthz} respon \texttt{\{ok:true\}}.
\end{itemize}

% ---------------------------------------------------------------------------
\section{Proves i diagnòstic ràpid}

\begin{itemize}[nosep]
  \item \textbf{Servidor}: \texttt{cd server \&\& node --test tests/*.test.js} (43
        proves: 13 de partida, 12 de recuperació, 5 de 2FA, 5 de mapa, 4 de política,
        3 de xarxa HTTP i 1 de seed).
  \item \textbf{Client}: \texttt{cd client \&\& node --test tests/*.test.js} (5 proves:
        2 de càmera, 2 de moviment i 1 de terreny).
  \item \textbf{Build}: \texttt{cd client \&\& npm run build}.
  \item \textbf{Tauler}: \texttt{npm run db} crea l'esquema i importa el seed si la
        taula és buida; no esborra dades.
\end{itemize}

\begin{table}[ht]
\small
\begin{tabularx}{\textwidth}{@{}p{1.2cm}X@{}}
\toprule
\textbf{Codi} & \textbf{Causa habitual} \\
\midrule
401 & Token absent, caducat, invàlid o repte de 2FA usat com a sessió. \\
403 & Compte bloquejat; el cos porta \texttt{banned}, \texttt{reason} i
\texttt{redirect}. \\
404 & La sala no existeix o el jugador ja no hi és. \\
409 & Sala plena, partida ja començada o 2FA ja actiu. \\
410 & El jugador ha estat expulsat de la sala (per exemple, per inactivitat). \\
429 & Límit de peticions o d'intents de 2FA. \\
400 & Camps invàlids, contrasenya curta, codi TOTP incorrecte o enllaç caducat. \\
\bottomrule
\end{tabularx}
\end{table}

\textbf{Errors típics del projecte:}

\begin{itemize}[nosep]
  \item El joc no es connecta: comprova que el servidor escolta al 2567 i que
        \texttt{/config.js} apunta a la URL correcta; en local, el client usa
        \texttt{location.hostname}.
  \item Els canvis del mapa no es veuen: reinicia el servidor (el tmj es llegeix una
        sola vegada) i mantingues \texttt{GRID = 60} al client.
  \item El jugador desapareix: ha estat 15~s sense sondejar; el \texttt{sendBeacon}
        només cobreix el tancament normal.
  \item Un 2FA no valida: el rellotge del dispositiu ha d'estar dins del període de
        30~s; el codi caduca i el repte de login dura 5~min.
  \item Un usuari banejat encara té una pàgina oberta: la propera petició a l'API rep
        403; la comprovació és a la base de dades, no només al token.
\end{itemize}

\end{document}
